\uuid{XU3R}
\exo7id{676}
\titre{exo7 676}
\auteur{cousquer}
\organisation{exo7}
\datecreate{2003-10-01}
\isIndication{false}
\isCorrection{false}
\chapitre{Continuité, limite et étude de fonctions réelles}
\sousChapitre{Continuité : pratique}
\cpge{mpsi-mp2i.fonctions-valeurs-reelles-complexes.a, pcsi.techniques-calcul-fonctions.a, ptsi.techniques-calcul-fonctions.a, mpsi-mp2i.fonctions-valeurs-reelles-complexes.b, pcsi.techniques-calcul-fonctions.b, ptsi.techniques-calcul-fonctions.b, mpsi-mp2i.limites-continuite.c, pcsi.limites-continuite.c, ptsi.limites-continuite.c, bcpst1.fonctions-reelles-usuelles.a, bcpst1.limites-continuite.c, bcpst1.limites-continuite.d, bcpst1.derivation.a}
\module{Analyse}
\niveau{L1}
\difficulte{}

\contenu{
\question{
Montrer que l'application $f : \mathbb{R} \rightarrow \mathbb{R}$ définie par
$\displaystyle{ f(x)=\frac{x}{1+\vert x\vert}}$
est strictement croissante puis que pour tout $y\in\quad]-1,1[ $ il existe
un unique $x\in \mathbb{R}$
tel que $f(x)=y$.
}
}
